Why does a bounded sequence of functions sometimes refuse to converge, and what survives when it does? Infinite dimensions
break the comfortable rules of linear algebra, and the repairs built between 1900 and 1935 are among the most useful ideas
in mathematics.

This book proves them. Hahn--Banach, Baire's three theorems and weak compactness; Hilbert spaces and the spectral theorem
for compact, bounded and unbounded operators; Fourier analysis from the circle to the real line, to distributions, to finite
groups and fast algorithms, and on to locally compact groups, Haar measure and the $p$-adic numbers. Then complex analysis
in earnest: the Jordan curve theorem, Cauchy's theorem in its final form, the growth of entire functions, the Riemann
mapping theorem, and the Mellin transform that measures the cost of algorithms. It ends with Tauberian theorems, which read
a sequence off its transform, and prove the prime number theorem again from almost nothing.

Every theorem is proved in full, every exercise has a hint and a reasoned solution. Written for the pleasure of seeing why things are true.
